% GENERATED FILE. Edit canonical lessons, then run npm run generate:books.
% canonical-source-hash: fnv1a64:c07fcbf3883cdbe0
% canonical-lessons: TE-C86-slowly, TE-C86-oncemore, TE-C86-loudly, TE-C86-oneminute, TE-C86-wait

\chapter{Slower, Please}
\label{ch:te-slower-please}
\hlchaptermodality{86}

\begin{chapteropening}
\textbf{By the end of this chapter:} \emph{I can ask someone to speak slowly or loudly, to say it once more, or to wait a minute.}

Complete TE-C86-wait: keep a conversation going when you cannot follow it.
\end{chapteropening}

\section[nemmadigā]{\te{నెమ్మదిగా} (nemmadigā) — slowly}

\label{lesson:TE-C86-slowly}



\begin{quote}
Before the new one: say the Telugu for anger, then the Telugu for sorrow.
\end{quote}

\subsection*{You'll want to know: \te{నెమ్మదిగా}}
\textbf{\te{నెమ్మదిగా}} (\emph{nemmadigā}) — \textquotedblleft{}slowly\textquotedblright{}.

\textbf{\te{నెమ్మదిగా} \te{మాట్లాడండి}} — \textquotedblleft{}please speak slowly\textquotedblright{}. You have the verb from your first chapters and the polite ending from \emph{raṇḍi}.

Notice the \textbf{-\te{గా}} at the end: the same ending that made \emph{alasaṭagā} and \emph{caligā}.

\begin{grammarlens}[title={What you've built}]
The request a beginner needs most.
\end{grammarlens}

\subsection*{Your turn}
\begin{itemize}
\raggedright
  \item \emph{Say it:} \emph{nemmadigā}
  \item \emph{Say it:} \emph{nemmadigā}, once more
  \item \emph{Say it:} say \emph{nemmadigā māṭlāḍaṇḍi}
  \item \emph{Recall:} say the Telugu for anger, then the Telugu for sorrow, then say \emph{nemmadigā} again
\end{itemize}

\begin{tcolorbox}[breakable,colback=teal!4,colframe=teal!35!black,title={Before you move on}]
What does \te{నెమ్మదిగా} mean? (\textquotedblleft{}Slowly\textquotedblright{}.) Say it once more.
\end{tcolorbox}
\section[iṅkōsāri]{\te{ఇంకోసారి} (iṅkōsāri) — once more}

\label{lesson:TE-C86-oncemore}



\begin{quote}
Before the new one: say the Telugu for sorrow, then the Telugu for slowly.
\end{quote}

\subsection*{You'll want to know: \te{ఇంకోసారి}}
\textbf{\te{ఇంకోసారి}} (\emph{iṅkōsāri}) — \textquotedblleft{}once more, one more time\textquotedblright{}.

\textbf{\te{ఇంకోసారి} \te{చెప్పండి}} — \textquotedblleft{}please say it once more\textquotedblright{}. \textbf{\te{సారి}} is \textquotedblleft{}a time, an occasion\textquotedblright{}, and \textbf{\te{ఇంకో}} is \textquotedblleft{}another\textquotedblright{}.

\begin{grammarlens}[title={What you've built}]
A way to hear it again without saying you did not understand.
\end{grammarlens}

\subsection*{Your turn}
\begin{itemize}
\raggedright
  \item \emph{Say it:} \emph{iṅkōsāri}
  \item \emph{Say it:} \emph{iṅkōsāri}, once more
  \item \emph{Say it:} say \emph{iṅkōsāri ceppaṇḍi}
  \item \emph{Recall:} say the Telugu for sorrow, then the Telugu for slowly, then say \emph{iṅkōsāri} again
\end{itemize}

\begin{tcolorbox}[breakable,colback=teal!4,colframe=teal!35!black,title={Before you move on}]
What does \te{ఇంకోసారి} mean? (\textquotedblleft{}Once more\textquotedblright{}.) Say it once more.
\end{tcolorbox}
\section[gaṭṭigā]{\te{గట్టిగా} (gaṭṭigā) — loudly; firmly}

\label{lesson:TE-C86-loudly}



\begin{quote}
Before the new one: say the Telugu for slowly, then the Telugu for once more.
\end{quote}

\subsection*{You'll want to know: \te{గట్టిగా}}
\textbf{\te{గట్టిగా}} (\emph{gaṭṭigā}) — \textquotedblleft{}loudly\textquotedblright{}, and also \textquotedblleft{}firmly, tightly\textquotedblright{}.

\textbf{\te{కొంచెం} \te{గట్టిగా} \te{మాట్లాడండి}} — \textquotedblleft{}please speak a little louder\textquotedblright{}. \textbf{\te{కొంచెం}} softens it, as it softens any request.

\begin{grammarlens}[title={What you've built}]
Slowly and loudly: the two fixes for a sentence you could not catch.
\end{grammarlens}

\subsection*{Your turn}
\begin{itemize}
\raggedright
  \item \emph{Say it:} \emph{gaṭṭigā}
  \item \emph{Say it:} \emph{gaṭṭigā}, once more
  \item \emph{Say it:} say \emph{koñcem gaṭṭigā māṭlāḍaṇḍi}
  \item \emph{Recall:} say the Telugu for slowly, then the Telugu for once more, then say \emph{gaṭṭigā} again
\end{itemize}

\begin{tcolorbox}[breakable,colback=teal!4,colframe=teal!35!black,title={Before you move on}]
What does \te{గట్టిగా} mean? (\textquotedblleft{}Loudly; firmly\textquotedblright{}.) Say it once more.
\end{tcolorbox}
\section[okka nimiṣaṁ]{\te{ఒక్క} \te{నిమిషం} (okka nimiṣaṁ) — one minute, a moment}

\label{lesson:TE-C86-oneminute}



\begin{quote}
Before the new one: say the Telugu for once more, then the Telugu for loudly.
\end{quote}

\subsection*{You'll want to know: \te{ఒక్క} \te{నిమిషం}}
\textbf{\te{ఒక్క} \te{నిమిషం}} (\emph{okka nimiṣaṁ}) — \textquotedblleft{}one minute\textquotedblright{}, said the way English says \textquotedblleft{}a moment\textquotedblright{}.

\textbf{\te{నిమిషం}} is \textquotedblleft{}a minute\textquotedblright{}, and \textbf{\te{ఒక్క}} is \textquotedblleft{}a single one\textquotedblright{}, a close relative of \textbf{\te{ఒకటి}}, the \textquotedblleft{}one\textquotedblright{} you count with.

\begin{grammarlens}[title={What you've built}]
A way to buy time while you find the word.
\end{grammarlens}

\subsection*{Your turn}
\begin{itemize}
\raggedright
  \item \emph{Say it:} \emph{okka nimiṣaṁ}
  \item \emph{Say it:} \emph{okka nimiṣaṁ}, once more
  \item \emph{Say it:} say \emph{okka nimiṣaṁ}
  \item \emph{Recall:} say the Telugu for once more, then the Telugu for loudly, then say \emph{okka nimiṣaṁ} again
\end{itemize}

\begin{tcolorbox}[breakable,colback=teal!4,colframe=teal!35!black,title={Before you move on}]
What does \te{ఒక్క} \te{నిమిషం} mean? (\textquotedblleft{}One minute, a moment\textquotedblright{}.) Say it once more.
\end{tcolorbox}
\section[āgaṇḍi]{\te{ఆగండి} (āgaṇḍi) — please wait}

\label{lesson:TE-C86-wait}



\begin{quote}
Before the new one: say the Telugu for loudly, then the Telugu for one minute.
\end{quote}

\subsection*{You'll want to know: \te{ఆగండి}}
\textbf{\te{ఆగండి}} (\emph{āgaṇḍi}) — \textquotedblleft{}please wait\textquotedblright{}, or \textquotedblleft{}please stop\textquotedblright{}.

It is the verb \textbf{\te{ఆగు}} with the polite \textbf{-\te{అండి}}. \textbf{\te{ఆగండి}, \te{ఒక్క} \te{నిమిషం}} — \textquotedblleft{}wait, one minute\textquotedblright{}.

\begin{grammarlens}[title={What you've built}]
Wait, slower, louder, once more: enough to hold a conversation you cannot yet follow.
\end{grammarlens}

\subsection*{Your turn}
\begin{itemize}
\raggedright
  \item \emph{Say it:} \emph{āgaṇḍi}
  \item \emph{Say it:} \emph{āgaṇḍi}, once more
  \item \emph{Say it:} say \emph{āgaṇḍi, okka nimiṣaṁ}
  \item \emph{Recall:} say the Telugu for loudly, then the Telugu for one minute, then say \emph{āgaṇḍi} again
\end{itemize}

\begin{tcolorbox}[breakable,colback=teal!4,colframe=teal!35!black,title={Before you move on}]
What does \te{ఆగండి} mean? (\textquotedblleft{}Please wait; please stop\textquotedblright{}.) Say it once more.
\end{tcolorbox}
